\documentclass[11pt]{scrreprt}
%%fakesection Load packages
% These are evan.sty
\usepackage{amsmath,amssymb,amsthm}
\usepackage{mathrsfs}
\usepackage[usenames,svgnames,dvipsnames]{xcolor}
\usepackage{hyperref}
\usepackage[nameinlink]{cleveref}
\usepackage{textcomp}
\usepackage{enumerate}
\usepackage[obeyFinal,textsize=scriptsize,shadow]{todonotes}
\usepackage{mathtools}
\usepackage{microtype}

\usepackage[normalem]{ulem}
\usepackage{stmaryrd}
\usepackage{wasysym}
\usepackage{multirow}
\usepackage{prerex}

%%fakesection evan.sty macros
%Small commands
\newcommand{\cbrt}[1]{\sqrt[3]{#1}}
\newcommand{\floor}[1]{\left\lfloor #1 \right\rfloor}
\newcommand{\ceiling}[1]{\left\lceil #1 \right\rceil}
\newcommand{\mailto}[1]{\href{mailto:#1}{\texttt{#1}}}
\newcommand{\ol}{\overline}
\newcommand{\ul}{\underline}
\newcommand{\wt}{\widetilde}
\newcommand{\wh}{\widehat}
\newcommand{\eps}{\varepsilon}
%\renewcommand{\iff}{\Leftrightarrow}
%\renewcommand{\implies}{\Rightarrow}
\newcommand{\vocab}[1]{\textbf{\color{blue} #1}}
\newcommand{\catname}{\mathsf}
\newcommand{\hrulebar}{
  \par\hspace{\fill}\rule{0.95\linewidth}{.7pt}\hspace{\fill}
  \par\nointerlineskip \vspace{\baselineskip}
}
\newcommand{\half}{\frac{1}{2}}

%For use in author command
\newcommand{\plusemail}[1]{\\ \normalfont \texttt{\mailto{#1}}}

%More commands and math operators
\DeclareMathOperator{\cis}{cis}
\DeclareMathOperator*{\lcm}{lcm}
\DeclareMathOperator*{\argmin}{arg min}
\DeclareMathOperator*{\argmax}{arg max}

%Convenient Environments
\newenvironment{soln}{\begin{proof}[Solution]}{\end{proof}}
\newenvironment{parlist}{\begin{inparaenum}[(i)]}{\end{inparaenum}}
\newenvironment{gobble}{\setbox\z@\vbox\bgroup}{\egroup}

%Inequalities
\newcommand{\cycsum}{\sum_{\mathrm{cyc}}}
\newcommand{\symsum}{\sum_{\mathrm{sym}}}
\newcommand{\cycprod}{\prod_{\mathrm{cyc}}}
\newcommand{\symprod}{\prod_{\mathrm{sym}}}

%From H113 "Introduction to Abstract Algebra" at UC Berkeley
\newcommand{\CC}{\mathbb C}
\newcommand{\FF}{\mathbb F}
\newcommand{\NN}{\mathbb N}
\newcommand{\QQ}{\mathbb Q}
\newcommand{\RR}{\mathbb R}
\newcommand{\ZZ}{\mathbb Z}
\newcommand{\charin}{\text{ char }}
\DeclareMathOperator{\sign}{sign}
\DeclareMathOperator{\Aut}{Aut}
\DeclareMathOperator{\Inn}{Inn}
\DeclareMathOperator{\Syl}{Syl}
\DeclareMathOperator{\Gal}{Gal}
\DeclareMathOperator{\GL}{GL} % General linear group
\DeclareMathOperator{\SL}{SL} % Special linear group

%From Kiran Kedlaya's "Geometry Unbound"
\newcommand{\abs}[1]{\left\lvert #1 \right\rvert}
\newcommand{\norm}[1]{\left\lVert #1 \right\rVert}
\newcommand{\dang}{\measuredangle} %% Directed angle
\newcommand{\ray}[1]{\overrightarrow{#1}} 
\newcommand{\seg}[1]{\overline{#1}}
\newcommand{\arc}[1]{\wideparen{#1}}

%From M275 "Topology" at SJSU
\newcommand{\id}{\mathrm{id}}
\newcommand{\taking}[1]{\xrightarrow{#1}}
\newcommand{\inv}{^{-1}}

%From M170 "Introduction to Graph Theory" at SJSU
\DeclareMathOperator{\diam}{diam}
\DeclareMathOperator{\ord}{ord}
\newcommand{\defeq}{\overset{\mathrm{def}}{=}}

%From the USAMO .tex filse
\newcommand{\st}{^{\text{st}}}
\newcommand{\nd}{^{\text{nd}}}
\renewcommand{\th}{^{\text{th}}}
\newcommand{\dg}{^\circ}
\newcommand{\ii}{\item}

% From Math 55 and Math 145 at Harvard
\newenvironment{subproof}[1][Proof]{%
\begin{proof}[#1] \renewcommand{\qedsymbol}{$\blacksquare$}}%
{\end{proof}}

\newcommand{\liff}{\leftrightarrow}
\newcommand{\lthen}{\rightarrow}
\newcommand{\opname}{\operatorname}
\newcommand{\surjto}{\twoheadrightarrow}
\newcommand{\injto}{\hookrightarrow}
\newcommand{\On}{\mathrm{On}} % ordinals
\DeclareMathOperator{\img}{im} % Image
\DeclareMathOperator{\Img}{Im} % Image
\DeclareMathOperator{\coker}{coker} % Cokernel
\DeclareMathOperator{\Coker}{Coker} % Cokernel
\DeclareMathOperator{\Ker}{Ker} % Kernel
\DeclareMathOperator{\rank}{rank}
\DeclareMathOperator{\Spec}{Spec} % spectrum
\DeclareMathOperator{\Tr}{Tr} % trace
\DeclareMathOperator{\pr}{pr} % projection
\DeclareMathOperator{\ext}{ext} % extension
\DeclareMathOperator{\pred}{pred} % predecessor
\DeclareMathOperator{\dom}{dom} % domain
\DeclareMathOperator{\ran}{ran} % range
\DeclareMathOperator{\Hom}{Hom} % homomorphism
\DeclareMathOperator{\End}{End} % endomorphism

% Things Lie
\newcommand{\kb}{\mathfrak b}
\newcommand{\kg}{\mathfrak g}
\newcommand{\kh}{\mathfrak h}
\newcommand{\kn}{\mathfrak n}
\newcommand{\ku}{\mathfrak u}
\newcommand{\kz}{\mathfrak z}
\DeclareMathOperator{\Ext}{Ext} % Ext functor
\DeclareMathOperator{\Tor}{Tor} % Tor functor
\newcommand{\gl}{\opname{\mathfrak{gl}}} % frak gl group
\renewcommand{\sl}{\opname{\mathfrak{sl}}} % frak gl group

% More script letters etc.
\newcommand{\SA}{\mathscr A}
\newcommand{\SB}{\mathscr B}
\newcommand{\SC}{\mathscr C}
\newcommand{\SD}{\mathscr D}
\newcommand{\SF}{\mathscr F}
\newcommand{\SG}{\mathscr G}
\newcommand{\SH}{\mathscr H}
\newcommand{\OO}{\mathcal O}

%% Napkin commands
\newcommand{\prototype}[1]{
	\emph{{\color{red} Prototypical example for this section:} #1} \par\medskip
}
\newenvironment{moral}{%
	\begin{mdframed}[linecolor=green!70!black]%
	\bfseries\color{green!70!black}}%
	{\end{mdframed}}

%%% END EXCERPT OF EVAN.STY
%%fakesection Commutative diagrams
\usepackage{diagrams}
\newarrow{Inj}C--->
\newarrow{Surj}----{>>}
\newarrow{Proj}----{>>}
\newarrow{Embed}>--->
\newarrow{Incl}C--->
\newarrow{Mapsto}|--->
\newarrow{Isom}<--->
\newarrow{Dotted}....>
\newarrow{Dashed}{}{dash}{}{dash}>
\newarrow{Line}-----
\usepackage{tikz-cd}
\usetikzlibrary{decorations.pathmorphing}
\tikzcdset{row sep/normal=3.14em,column sep/normal=3.14em}

%%fakesection Page layout
\usepackage[headsepline]{scrlayer-scrpage}
\renewcommand{\headfont}{}
\addtolength{\textheight}{3.14cm}
\setlength{\footskip}{0.5in}
\setlength{\headsep}{10pt}
\ihead{Evan Chen}
\automark{section}
\chead{}
\ohead{\headmark}
\cfoot{\pagemark}

%%fakesection Fancy section and chapter heads
\addtokomafont{chapterprefix}{\raggedleft}
\renewcommand*{\sectionformat}{\color{purple}\S\thesection\autodot\enskip}

\RedeclareSectionCommand[beforeskip=0.5em]{chapter}
\renewcommand*{\chapterformat}{%
\mbox{\scalebox{1.5}{\chapappifchapterprefix{\nobreakspace}}%
\scalebox{2.718}{\color{purple}\thechapter\autodot}\enskip}}

\addtokomafont{partprefix}{\rmfamily}
\renewcommand*{\partformat}{\color{purple}\scalebox{2.5}{\thepart}%
	\enlargethispage{2em}}
\newcommand{\problemhead}{Problems to think about}

%%fakesection Mini ToC
\usepackage[tight]{minitoc}
\noptcrule
\mtcsetfont{parttoc}{section}{\small\rmfamily\upshape\mdseries}
\renewcommand*{\partheadstartvskip}{}
\renewcommand*{\partheadendvskip}{}
\renewcommand\beforeparttoc{\hspace{\fill}\rule{0.95\linewidth}{2pt}\hspace{\fill}}
\let\oldpart\part
\renewcommand{\part}[1]{\oldpart{#1}\parttoc\cleardoublepage}

%%fakesection SNSD
\newcommand{\snsd}[2][]%
{\includegraphics[#1]{/home/evan/Pictures/TopologicalGG/#2}}

%%fakesection Theorems
\usepackage{thmtools}
\usepackage[framemethod=TikZ]{mdframed}

\theoremstyle{definition}
\mdfdefinestyle{mdbluebox}{%
	roundcorner = 10pt,
	linewidth=1pt,
	skipabove=12pt,
	innerbottommargin=9pt,
	skipbelow=2pt,
	nobreak=true,
	linecolor=blue,
	backgroundcolor=TealBlue!5,
}
\declaretheoremstyle[
	headfont=\sffamily\bfseries\color{MidnightBlue},
	mdframed={style=mdbluebox},
	headpunct={\\[3pt]},
	postheadspace={0pt}
]{thmbluebox}

\mdfdefinestyle{mdrecbox}{%
	linewidth=0.5pt,
	skipabove=12pt,
	frametitleaboveskip=5pt,
	frametitlebelowskip=0pt,
	skipbelow=2pt,
	frametitlefont=\bfseries,
	innertopmargin=4pt,
	innerbottommargin=8pt,
	nobreak=true,
	linecolor=RawSienna,
	backgroundcolor=Salmon!5,
}
\declaretheoremstyle[
	headfont=\bfseries\color{RawSienna},
	mdframed={style=mdrecbox},
	headpunct={\\[3pt]},
	postheadspace={0pt},
]{thmrecbox}

\declaretheorem[%
style=thmbluebox,name=Theorem,numberwithin=section]{theorem}
\declaretheorem[style=thmbluebox,name=Lemma,sibling=theorem]{lemma}
\declaretheorem[style=thmbluebox,name=Proposition,sibling=theorem]{proposition}
\declaretheorem[style=thmbluebox,name=Corollary,sibling=theorem]{corollary}
\declaretheorem[style=thmrecbox,name=Example,sibling=theorem]{example}

\mdfdefinestyle{mdinlinebox}{%
	skipabove=8pt,
	linewidth=3pt,
	rightline=false,
	leftline=true,
	topline=false,
	bottomline=false,
	linecolor=gray,
	backgroundcolor=black!2
}
\declaretheoremstyle[
	headfont=\bfseries,
	bodyfont=\normalfont\small,
	spaceabove=0pt,
	spacebelow=0pt,
	mdframed={style=mdinlinebox}
]{thminlinebox}
\theoremstyle{theorem}
\declaretheorem[name=Question,sibling=theorem,style=thminlinebox]{ques}
\declaretheorem[name=Exercise,sibling=theorem,style=thminlinebox]{exercise}

\theoremstyle{definition}
\newtheorem{claim}[theorem]{Claim}
\newtheorem{definition}[theorem]{Definition}
\newtheorem{fact}[theorem]{Fact}
\newtheorem{remark}[theorem]{Remark}
\newtheorem{abuse}[theorem]{Abuse of Notation}

\newtheorem{problem}{Problem}[chapter]
\renewcommand{\theproblem}{\thechapter\Alph{problem}}
\newtheorem{sproblem}[problem]{Problem}
\newtheorem{dproblem}[problem]{Problem}
\renewcommand{\thesproblem}{\theproblem$^{\star}$}
\renewcommand{\thedproblem}{\theproblem$^{\dagger}$}
\newcommand{\listhack}{$\empty$\vspace{-2em}}

%%fakesection Answers
\usepackage{answers}
\Newassociation{hint}{answeritem}{tex/backmatter/all-hints}
\Newassociation{sol}{answeritem}{tex/backmatter/all-solns}
\renewcommand{\solutionextension}{out}
\renewenvironment{answeritem}[1]{\item[\bf #1.]}{}

%%fakesection Hot chili peppers
\reversemarginpar
\newcommand{\prechili}{\vspace*{0.3em}\hspace*{1.5em}}
\newcommand{\nochili}{\hspace*{1.5em}}
\newcommand{\chili}{\includegraphics[width=1.5em]{media/chili.png}}
\newcommand{\gim}{\marginpar{\prechili\nochili\nochili\chili}}
\newcommand{\yod}{\marginpar{\prechili\nochili\chili\chili}}
\newcommand{\kurumi}{\marginpar{\prechili\chili\chili\chili}}

%%fakesection Table of contents
% First add ToC to ToC
\makeatletter
\usepackage{etoolbox}
\pretocmd{\tableofcontents}{%
  \if@openright\cleardoublepage\else\clearpage\fi
  \pdfbookmark[0]{\contentsname}{toc}%
}{}{}%
\makeatother
\setcounter{tocdepth}{1}
% Fix spacing issues in To
\usepackage[tocindentauto]{tocstyle}
\usetocstyle{KOMAlike}

%%fakesection Asymptote definitions
\usepackage{patch-asy}
\numberwithin{asy}{chapter}
\renewcommand{\theasy}{\thechapter\Alph{asy}}
\begin{asydef}
	import olympiad;
	import cse5;
	size(6cm);
	usepackage("amsmath");
	usepackage("amssymb");
	defaultpen(fontsize(11pt));
	pointfontsize=11;
	void bigbox(string L) {
		filldraw((-4,3)--(4,3)--(4,-3)--(-4,-3)--cycle,
			lightblue+opacity(0.2), black);
		label(L, (4,3), dir(-45));
	}
	void bigblob(string L) {
		filldraw((-4,3)..(-5,1)..(-4.7,0)..(-4.1,-2)..(-4,-3)
			..(-1,-3.3)..(1,-3.2)..(4,-3)
			..(4.3,-1)..(4.1,0.5)..(4,3)
			..(2,3.3)..(0,3.1)..(-3,2.9)..cycle,
			lightblue+opacity(0.2), black);
		label(L, (4,3), dir(5));
	}

	real[] default_ticks = {-4,-3,-2,-1,0,1,2.3,4};
	void cplane(real theta=-45, real[] T = default_ticks, real xmin = -4.5, real xmax = 4.5, real S = 3) {
		graph.xaxis("Re", xmin, xmax, graph.Ticks(format="%", Ticks=T, Size=S), Arrows);
		graph.yaxis("Im", xmin, xmax, graph.Ticks(format="%", Ticks=T, Size=S), Arrows);
		Drawing("0", origin, dir(theta));
		markscalefactor *= xmax;
	}

	pair opendot(pair P, pen p = defaultpen) {
		dot(P, p, UnFill);
		dot(P, p, Draw);
		return P;
	}
\end{asydef}
\def\asydir{asy}

%%fakesection Bibliography
\usepackage[backend=biber,style=alphabetic]{biblatex}
\DeclareLabelalphaTemplate{
	\labelelement{
		\field[final]{shorthand}
		\field{label}
		\field[strwidth=2,strside=left]{labelname}
	}
	\labelelement{
		\field[strwidth=2,strside=right]{year}    
	}
}
\DeclareFieldFormat{labelalpha}{\textbf{\scriptsize #1}}
\addbibresource{references.bib}
\addbibresource{images.bib}

%%fakesection Napkin Macros
\newcommand{\Zc}[1]{\ZZ_{#1}}
\newcommand{\Zm}[1]{\ZZ_{#1}^\times}
\newcommand{\pre}{^{\text{pre}}}
\newcommand{\normalin}{\trianglelefteq}
\newcommand{\triv}{\mathrm{triv}}
\newcommand{\largeotimes}[1]{\mathop{\otimes}\limits_{#1}}
\newcommand{\Mat}{\mathrm{Mat}}
\newcommand{\PGL}[1]{\mathbf{PGL}_{#1}(\CC)}

\DeclareMathOperator{\Stab}{Stab}
\DeclareMathOperator{\FixPt}{FixPt}
\DeclareMathOperator{\refl}{refl}
\DeclareMathOperator{\Fun}{Fun}
\DeclareMathOperator{\Irrep}{Irrep}
\DeclareMathOperator{\Res}{Res}
\DeclareMathOperator{\Reg}{Reg}
\DeclareMathOperator{\Classes}{Classes}
\newcommand{\FunCl}{\Fun_{\mathrm{class}}}
\newcommand{\Homrep}{\Hom_{\mathrm{rep}}}
\newcommand{\ab}{^\text{ab}} % abelianization
\DeclareMathOperator{\ev}{ev}
\DeclareMathOperator{\Wind}{\mathbf I}
\newcommand{\Ctriv}{\CC_{\mathrm{triv}}}
\newcommand{\Csign}{\CC_{\mathrm{sign}}}

% Alg geom macros
\newcommand{\VV}{\mathcal V}
\newcommand{\Vp}{\mathcal V_{\text{pr}}}
\newcommand{\Aff}{\mathbb A}
\newcommand{\II}{\mathcal I_{\text{rad}}}
\newcommand{\RP}{\mathbb{RP}}
\newcommand{\CP}{\mathbb{CP}}
\newcommand{\Cl}{\mathrm{Cl}}
\newcommand{\restrict}[1]{\bgroup\restriction_{#1}\egroup}
\DeclareMathOperator{\Opens}{OpenSets}
\DeclareMathOperator{\Proj}{Proj}
\DeclareMathOperator{\res}{res}
\newcommand{\mm}{\mathfrak m}
\newcommand{\sh}{^\mathrm{sh}}

\renewcommand{\Re}{\opname{Re}}
\renewcommand{\Im}{\opname{Im}}

%% Category Theory macros
\renewcommand{\AA}{\mathcal A}
\newcommand{\BB}{\mathcal B}
\newcommand{\obj}{\operatorname{obj}}
\newcommand{\op}{^{\mathrm{op}}}


\usepackage{tkz-graph}
\pgfarrowsdeclare{biggertip}{biggertip}{%
  \setlength{\arrowsize}{1pt}  
  \addtolength{\arrowsize}{.1\pgflinewidth}  
  \pgfarrowsrightextend{0}  
  \pgfarrowsleftextend{-5\arrowsize}  
}{%
  \setlength{\arrowsize}{1pt}  
  \addtolength{\arrowsize}{.1\pgflinewidth}  
  \pgfpathmoveto{\pgfpoint{-5\arrowsize}{4\arrowsize}}  
  \pgfpathlineto{\pgfpointorigin}  
  \pgfpathlineto{\pgfpoint{-5\arrowsize}{-4\arrowsize}}  
  \pgfusepathqstroke  
}  
\tikzset{
	EdgeStyle/.style = {>=biggertip}
}

\usepackage[all,cmtip,2cell]{xy}
\UseTwocells
\newcommand\nattfm[5]{\xymatrix@C+2pc{#1 \rtwocell<4>^{#2}_{#4}{\; #3} & #5}}

%% Alg Top macros
\DeclareMathOperator{\Cells}{Cells}
\newcommand{\HdR}{H_{\mathrm{dR}}}

%% Alg NT macros
\renewcommand{\aa}{\mathfrak a}
\newcommand{\bb}{\mathfrak b}
\newcommand{\pp}{\mathfrak p}
\newcommand{\qq}{\mathfrak q}

\newcommand{\Frob}{\mathrm{Frob}}
\DeclareMathOperator{\Norm}{N}
\DeclareMathOperator{\Ram}{Ram}
\newcommand{\TrK}{\Tr_{K/\QQ}}
\newcommand{\NK}{\Norm_{K/\QQ}}
\newcommand{\kf}{\mathfrak f}
\newcommand{\kP}{\mathfrak P}
\newcommand{\kQ}{\mathfrak Q}

%% Diff geo macros
\newcommand{\ee}{\mathbf e}
\newcommand{\fpartial}[2]{\frac{\partial #1}{\partial #2}}

%% ST macros
\newcommand{\CH}{\mathsf{CH}}
\newcommand{\ZFC}{\mathsf{ZFC}}

\newcommand{\Name}{\text{Name}}
\newcommand{\Po}{\mathbb P}
\newcommand{\nrank}{\opname{n-rank}} % ranks
\newcommand{\PP}{\mathcal P}

\newcommand{\EmptySet}{\mathrm{EmptySet}}
\newcommand{\PowerSet}{\mathrm{PowerSet}}
\newcommand{\Pairing}{\mathrm{Pairing}}
\newcommand{\Infinity}{\mathrm{Infinity}}
\newcommand{\Extensionality}{\mathrm{Extensionality}}
\newcommand{\Foundation}{\mathrm{Foundation}}
\newcommand{\Union}{\mathrm{Union}}
\newcommand{\Comprehension}{\mathrm{Comprehension}}
\newcommand{\Replacement}{\mathrm{Replacement}}

\usepackage{mathrsfs}

\newcommand\MM{\mathscr M}
\newcommand\llex{<_{\text{lex}}}

\DeclareMathOperator{\cof}{cof}

%% Quantum macros
\usepackage{braket}
\newcommand{\cvec}[1]{\begin{bmatrix} #1 \end{bmatrix}}
\newcommand{\pair}[2]{\begin{bmatrix} #1 \\ #2 \end{bmatrix}}
\newcommand{\zup}{\ket\uparrow}
\newcommand{\zdown}{\ket\downarrow}
\newcommand{\xup}{\ket\rightarrow}
\newcommand{\xdown}{\ket\leftarrow}
\newcommand{\yup}{\ket\otimes}
\newcommand{\ydown}{\ket\odot}
\newcommand{\UCNOT}{U_{\mathrm{CNOT}}}
\newcommand{\UQFT}{U_{\mathrm{QFT}}}

\addbibresource{../../references.bib}

\newcommand{\CH}{\mathsf{CH}}
\newcommand{\ZFC}{\mathsf{ZFC}}

\newcommand{\Name}{\text{Name}}
\newcommand{\Po}{\mathbb P}
\newcommand{\nrank}{\opname{n-rank}} % ranks

\newcommand{\EmptySet}{\mathtt{EmptySet}}
\newcommand{\PowerSet}{\mathtt{PowerSet}}
\newcommand{\Extensionality}{\mathtt{Extensionality}}
\newcommand{\Foundation}{\mathtt{Foundation}}

\usepackage{mathrsfs}
\newcommand\MM{\mathscr M}

\def\asydir{}

\begin{document}
\title{Set Theory}
\maketitle

\chapter{ZFC and Ordinal Numbers}
\section{The Ultimate Functional Equation}
definition nesting. . .

\section{The Axioms}

\section{Sets vs Classes}

\section{Ordinal Numbers}
Define
spiral picture

Induction

\section\problemhead
\begin{problem}
	Show that the class of ordinals $\On$ is a proper class.
\end{problem}


\chapter{Ordinals and Cardinals}
transfinite induction

ordinal arithmetic

definition of cardinals

(note that they depend on the outside universe!)

cardinal arithmetic

cofinality: regular vs singular

true to its name, regular is the normal variety

\section\problemhead
Inaccessible


\chapter{Inner Model Theory}
Model theory is \emph{really} meta, so you will have to pay attention here.

Roughly, a ``model of $\ZFC$'' is a set with a binary relation that satisfies the $\ZFC$ axioms,
just as a group is a set with a binary operation that satisfies the group axioms.
Unfortunately, unlike with groups, it is very hard for me to give interesting examples of models,
for the simple reason that we are literally trying to model the entire universe.

\section{Model}
\prototype{$(\omega, \in)$ obeys $\PowerSet$, $V_\kappa$ is a model for $\kappa$ inaccessible (later).}
\begin{definition}
	A \vocab{model} $\MM$ consists of a set $M$ and a
	binary relation $E \subseteq M \times M$.
	(The $E$ relation is the ``$\in$'' for the model.)
\end{definition}
\begin{remark}
	I'm only considering \emph{set-sized} models where $M$ is a set.
	Experts may be aware that I can actually play with $M$ being a class,
	but that would require too much care for now.
\end{remark}
If you have a model, you can ask certain things about it.
For example, you can ask ``does it satisfy $\EmptySet$?''.
Let me give you an example of what I mean, and then make it rigorous.
\begin{example}
	[A Stupid Model]
	Let's take $\MM = (M,E) = \left( \omega, \in \right)$.
	This is not a very good model of $\ZFC$, but let's see if we can
	make sense of some of the first few axioms.
	\begin{enumerate}[(a)]
		\ii $\MM$ satisfies $\mathtt{Extensionality}$, which is the sentence
		\[ \forall x \forall y \forall a : \left( a \in x \iff a \in y \right) \implies x = y. \]
		This just follows from the fact that $E$ is actually $\in$.

		\ii $\MM$ satisfies $\EmptySet$, which is the sentence
		\[ \exists a : \forall x \; \neg (x \in a). \]
		Namely, take $a = \varnothing \in \omega$.

		\ii $\MM$ does not satisfy $\mathtt{Pairing}$, since $\{1,3\}$ is not in $\omega$,
		even though $1, 3 \in \omega$.

		\ii Miraculously, $\MM$ satisfies $\mathtt{Union}$, since for any $n \in \omega$,
		$\cup n$ is $n-1$ (unless $n=0$).
		The Union axiom statements that
		\[ \forall a \exists z \quad \forall x \; (x \in z) \iff (\exists y : x \in y \in z). \]
		An important thing to notice is that the ``$\forall a$'' ranges only over
		the sets in the model of the universe, $\MM$.
	\end{enumerate}
\end{example}
\begin{example}
	[Very Important: This Stupid Model Satisfies $\PowerSet$]
	Most incredibly of all: $\MM = (\omega, \in)$ satisfies $\PowerSet$.
	This is a really important example.
	
	You might think this is ridiculous. Look at $2 = \{0,1\}$.
	The power set of this is $\{0, 1, 2, \{1\}\}$ which is not in the model, right?

	Well, let's look more closely at $\PowerSet$. It states that:
	\[ \forall x \exists a \forall y (y \in a \iff y \subseteq x). \]
	What happens if we set $x = 2 = \{0,1\}$?
	Well, actually, we claim that $a = 3 = \{0,1,2\}$ works.
	The key point is ``for all $y$'' -- this \emph{only ranges over the objects in $\MM$}.
	In $\MM$, the only subsets of $2$ are $0 = \varnothing$,
	$1 = \{0\}$ and $2 = \{0,1\}$.
	The ``set'' $\{1\}$ in the ``real world'' (in $V$) is not a set in the model $\MM$.

	In particular, you might say that in this strange new world,
	we have $2^n = n+1$, since $n = \{0,1,\dots,n-1\}$ really does
	have only $n+1$ subsets.
\end{example}

\begin{example}
	[Sentences with Parameters]
	The sentences we ask of our model are allowed to have ``parameters'' as well.
	For example, if $\MM = (\omega, \in)$ as before then $\MM$ satisfies the sentence
	\[ \forall x \in 3 (x \in 5). \]
\end{example}

\section{Sentences and Satisfaction}
With this intuitive notion, we can define what it means for a model to satisfy a sentence.
\begin{definition}
Note that any sentence $\phi$ can be written in one of the following five forms:
\begin{itemize}
	\ii $x \in y$
	\ii $x = y$
	\ii $\neg \psi$ (``not $\psi$'') for some shorter sentence $\psi$
	\ii $\psi_1 \lor \psi_2$ (```$\psi_1$ or $\psi_2$'')
	for some shorter sentences $\psi_1$, $\psi_1$
	\ii $\exists x \psi$ (``exists $x$'') for some shorter sentence $\psi$.
\end{itemize}
\end{definition}
\begin{ques}
	What happened to $\land$ (and) and $\forall$ (for all)?
	(Hint: use $\neg$.)
\end{ques}
Often (almost always, actually) we will proceed by so-called ``induction on formula complexity'',
meaning that we define or prove something by induction using this.
Note that we require all formulas to be finite.

Now suppose we have a sentence $\phi$, like $a = b$ or $\exists a \forall x \neg (x \in a)$,
plus a model $\MM = (M,E)$.
We want to ask whether $\MM$ satisfies $\phi$.

To give meaning to this, we have to designate certain variables as \vocab{parameters}.
For example, if I asked you 
\begin{quote}
	``Does $a=b$?''
\end{quote}
the first question you would ask is what $a$ and $b$ are.
So $a$, $b$ would be parameters: I have to give them values for this sentence to make sense.

On the other hand, if I asked you
\begin{quote}
	``Does $\exists a \forall x \neg (x \in a)$?''
\end{quote}
then you would just say ``yes''.
In this case, $x$ and $a$ are \emph{not} parameters.
In general, parameters are those variables whose meaning is not given by some $\forall$ or $\exists$.

In what follows, we will let $\phi(x_1, \dots, x_n)$ denote a formula $\phi$,
whose parameters are $x_1$, \dots, $x_n$.
Note that possibly $n=0$, for example all $\ZFC$ axioms have no parameters.

\begin{ques}
	Try to guess the definition of satisfaction before reading it below.
	(It's not very hard to guess!)
\end{ques}

\begin{definition}
	Let $\MM=(M,E)$ be a model.
	Let $\phi(x_1, \dots, x_n)$ be a sentence, and let $b_1, \dots, b_n \in M$.
	We will define a relation
	\[ \MM \vDash \phi[b_1, \dots, b_n] \]
	and say $\MM$ \vocab{satisfies} the sentence $\phi$ with parameters $b_1, \dots, b_n$.

	The relationship is defined by induction on formula complexity as follows:
	\begin{itemize}
		\ii If $\phi$ is ``$x_1=x_2$'' then $\MM \vDash \phi[b_1, b_2] \iff b_1 = b_2$.
		\ii If $\phi$ is ``$x_1\in x_2$'' then $\MM \vDash \phi[b_1, b_2] \iff b_1 \; E \; b_2$. \\
		(This is what we mean by ``$E$ interprets $\in$''.)
		\ii If $\phi$ is ``$\neg \psi$'' then
		$\MM \vDash \phi[b_1, \dots, b_n] \iff \MM \not\vDash \phi[b_1, \dots, b_n]$.
		\ii If $\phi$ is ``$\psi_1 \lor \psi_2$'' then $\MM \vDash \phi[b_1, \dots, b_n]$
		means $\MM \vDash \psi_i[b_1, \dots, b_n]$ for some $i=1,2$.
		\ii Most important case: suppose $\phi$ is $\exists x \psi(x,x_1, \dots, x_n)$.
		Then $\MM \vDash \phi[b_1, \dots, b_n]$ if and only if 
		\[ \exists b \in M \text{ such that } \MM \vDash \psi[b, b_1, \dots, b_n]. \]
		Note that $\psi$ has one extra parameter.
	\end{itemize}
\end{definition}
Notice where the information of the model actually gets used.
We only ever use $E$ in interpreting $x_1 \in x_2$; unsurprising.
But we only ever use the set $M$ when we are running over $\exists$ (and hence $\forall$).
That's well-worth keeping in mind:
\begin{moral}
	The behavior of a model essentially comes from $\exists$ and $\forall$,
	which search through the entire model $M$.
\end{moral}

And finally,
\begin{definition}
	A \vocab{model of $\ZFC$} is a model $\MM = (M,E)$ satisfying all $\ZFC$ axioms.
\end{definition}

We are especially interested in models of the form $(M, \in)$, where $M$ is a \emph{transitive} set.
(We want our universe to be transitive,-
otherwise we would have elements of sets which are not themselves
in the universe, which is very strange.)
Such a model is called a \vocab{transitive model}.
\begin{abuse}
	If $M$ is a transitive set, the model $(M, \in)$ will be abbreviated to just $M$.
\end{abuse}
\begin{definition}
	An \vocab{inner model} of $\ZFC$ is a transitive model satisfying $\ZFC$.
\end{definition}

\section{The Levy Hierarchy}
\prototype{$\mathtt{isSubset}(x,y)$ is absolute. The axiom
$\EmptySet$ is $\Sigma_1$, $\mathtt{isPowerSetOf}(X,x)$ is $\Pi_1$.}
A key point to remember is that the behavior of a model is largely determined by $\exists$ and $\forall$.
It turns out we can say even more than this.

Consider a formula such as
\[ \mathtt{isempty}(x) : \neg \exists a (a \in x) \]
which checks whether a given set $x$ has a nonempty element.
Technically, this has an ``$\exists$'' in it.
But somehow this $\exists$ does not really search over the entire model,
because it is \emph{bounded} to search in $x$.
That is, we might informally rewrite this as
\[ \neg (\exists x \in a) \]
which doesn't fit into the strict form,
but points out that we are only looking over $a \in x$.
We call such a quantifier a \vocab{bounded quantifier}.
%and write it
%in the way we see it in most mathematics, such as
%\[ (\exists x \in a) \quad\text{or}\quad (\forall x \in a). \]
%To be painfully explicit,
%$\exists x \in a \psi$ is short for $\exists x (x \in a \land \psi)$,
%while $\forall x \in a \psi$ is short for $\forall x (x \in a \implies \psi)$.

We like sentences with bounded quantifiers because they designate
properties which are \vocab{absolute} over transitive models.
It doesn't matter how strange your surrounding model $M$ is.
As long as $M$ is transitive, 
\[ M \vDash \mathtt{isEmpty}[\varnothing] \]
will always hold.
Similarly, the sentence
\[ \mathtt{isSubset}(x,y) : x \subseteq y \text { i.e. } \forall a \in x (a \in y). \]
Sentences with this property are called $\Sigma_0$ or $\Pi_0$.

The situation is different with a sentence like
\[
	\mathtt{isPowerSetOf}(y,x) :
	\forall z \left( z \subseteq x \iff z \in y  \right)
\]
which in English means ``$y$ is the power set of $x$'', or just $y = \PP(x)$.
The $\forall z$ is \emph{not} bounded here.
This weirdness is what allows things like
\[ \omega \vDash \text{``$\{0,1,2\}$ is the power set of $\{0,1\}$''} \]
and hence
\[ \omega \vDash \PowerSet \]
which was our stupid example earlier.
The sentence $\mathtt{isPowerSetOf}$ consists of an unbounded $\forall$ followed
by an absolute sentence, so we say it is $\Pi_1$.

More generally, the \vocab{Levy hierarchy} keeps track of how bounded our
quantifiers are.
Specifically,
\begin{itemize}
	\ii Formulas which have only bounded quantifiers are $\Delta_0 = \Sigma_0 = \Pi_0$.
	\ii Formulas of the form $\exists x_1 \dots \exists x_k \psi$ where $\psi$ is $\Pi_n$
	are consider $\Sigma_{n+1}$.
	\ii Formulas of the form $\forall x_1 \dots \forall x_k \psi$ where $\psi$ is $\Sigma_n$
	are consider $\Pi_{n+1}$.
\end{itemize}
(A formula which is both $\Sigma_n$ and $\Pi_n$ is said $\Delta_n$, but we won't
use this except for $n=0$.)

\begin{example}
	[Examples of $\Delta_0$ Sentences]
	\listhack
	\begin{enumerate}[(a)]
		\ii The sentences $\mathtt{isEmpty}(x)$, $x \subseteq y$, as discussed above.
		\ii The formula ``$x$ is transitive'' can be expanded as a $\Delta_0$ sentence.
		\ii The formula ``$x$ is an ordinal'' can be expanded as a $\Delta_0$ sentence.
	\end{enumerate}
\end{example}
\begin{exercise}
	Write out the expansions for ``$x$ is transitive'' and ``$x$ is ordinal''
	in a $\Delta_0$ form.
\end{exercise}
\begin{example}
	[More Complex Formulas]
	\listhack
	\begin{enumerate}[(a)]
		\ii The axiom $\EmptySet$ is $\Sigma_1$; it is $\exists a (\mathtt{isEmpty}(a))$,
		and $\mathtt{isEmpty}(a)$ is $\Delta_0$.
		\ii The formula ``$y = \PP(x)$'' is $\Pi_1$, as discussed above.
		\ii The formula ``$x$ is countable'' is $\Sigma_1$.
		One way to phrase it is ``$\exists f$ an injective map $x \injto \omega$'',
		which necessarily has an unbounded ``$\exists f$''.
		\ii The axiom $\PowerSet$ is $\Pi_3$:
		\[ \forall y \exists P \forall x (x\subseteq y \iff x \in P). \]
	\end{enumerate}
\end{example}

\section{Substructures, and Tarski-Vaught}
Let $\MM_1 = (M_1, E_1)$ and $\MM_2 = (M_2, E_2)$ be models.
\begin{definition}
	We say that $\MM_1 \subseteq \MM_2$ if $M_1 \subseteq M_2$ and
	$E_1$ agrees with $E_2$; we say $\MM_1$ is a \vocab{substructure} of $\MM_2$.
\end{definition}

That's boring. The good part is:
\begin{definition}
	We say $\MM_1 \prec \MM_2$, or $\MM_1$ is an \vocab{elementary substructure} of $\MM_2$,
	if for \emph{every} sentence $\phi(x_1, \dots, x_n)$
	and parameters $b_1, \dots, b_n \in M_1$, we have
	\[
		\MM_1 \vDash \phi[b_1, \dots, b_n] \iff
		\MM_2 \vDash \phi[b_1, \dots, b_n].
	\]
\end{definition}
In other words, $\MM_1$ and $\MM_2$ agree on every sentence possible.
Note that the $b_i$ have to come from $M_1$; if the $b_i$ came from $\MM_2$ then
asking something of $\MM_1$ wouldn't make sense.

Let's ask now: how would $\MM_1 \prec \MM_2$ fail to be true?
If we look at the possibly sentences, none of the atomic formulas,
nor the ``$\land$'' and ``$\neg$'', are going to cause issues.

The intuition you should be getting by now is that things go
wrong once we hit $\forall$ and $\exists$.
They won't go wrong for bounded quantifiers.
But unbounded quantifiers search the entire model, and that's where things go wrong.

To give a ``concrete example'':
imagine $\MM_1$ is MIT, and $\MM_2$ is the state of Massachusetts.
If $\MM_1$ thinks there exist hackers at MIT,
certainly there exist hackers in Massachusetts.
Where things go wrong is something like:
\[ \MM_2 \vDash \text{``$\exists x$ : $x$ is a course numbered $> 50$.''}. \]
This is true for $\MM_2$ because we can take the witness $x = \text{Math 55}$, say.
But it's false for $\MM_1$, because at MIT all courses are numbered $18.701$ or something similar.
\begin{moral}
	The issue is that the \emph{witness}
	for statements in $\MM_2$ do not necessarily propagate up
	down to witnesses for $\MM_1$, even though they do from $\MM_1$ to $\MM_2$.
\end{moral}

The Tarski-Vaught test says this is the only impediment:
if every witness in $\MM_2$ can be replaced by one in $\MM_1$ then $\MM_1 \prec \MM_2$.
\begin{lemma}
	[Tarski-Vaught]
	Let $\MM_1 \subseteq \MM_2$.
	Then $\MM_1 \prec \MM_2$ if and only if for
	every sentence $\phi(x, x_1, \dots, x_n)$ and parameters $b_1, \dots, b_n \in M_1$:
	if there is a witness $\tilde b \in M_2$ to $\MM_2 \vDash \phi(\tilde b, b_1 \dots, b_n)$
	then there is a witness $b \in M_1$ to $\MM_1 \vDash \phi(b, b_1, \dots, b_n)$.
\end{lemma}
\begin{proof}
	Easy after the above discussion.
	To formalize it, use induction on formula complexity.
\end{proof}

\section{Obtaining the Axioms of $\ZFC$}
Extending the above ideas, one can obtain without much difficulty the following.
The idea is that almost all the $\ZFC$ axioms are just $\Sigma_1$ claims about certain desired sets,
and so verifying an axiom reduces to checking some appropriate ``closure'' condition:
that the witness to the axiom is actually in the model.

For example, the $\EmptySet$ axiom is ``$\exists a (\mathtt{isEmpty}(a))$'',
and so we're happy as long as $\varnothing \in M$, which is of course
true for any nonempty transitive set $M$.

\begin{lemma}[Transitive Sets Inheriting $\ZFC$]
	\label{lem:transitive_ZFC}
	Let $M$ be a nonempty transitive set. Then
	\begin{enumerate}[(i)]
		\ii $M$ satisfies $\Extensionality$, $\Foundation$, $\EmptySet$.
		\ii $M \vDash \mathtt{Pairing}$ if $x,y \in M \implies \{x,y\} \in M$.
		\ii $M \vDash \mathtt{Union}$ if $x \in M \implies \cup x \in M$.
		\ii $M \vDash \PowerSet$ if $x \in M \implies \PP(x) \cap M \in M$.
		\ii $M \vDash \mathtt{Comprehension}$ if for every set $C$
		which is $M$-definable with parameters, and every $x \in M$, we have $C \cap x \in M$.
		\ii $M \vDash \mathtt{Collection}$ if for every function $F : M \to M$
		which is $M$-definable with parameters,
		and every $x \in M$, we have $F``x \in M$ as well.
		\ii $M \vDash \mathtt{Infinity}$ as long as $\omega \in M$.
	\end{enumerate}
\end{lemma}
Here, a set $X \subseteq M$ is \vocab{$M$-definable with parameters}
if it can be realized as
\[ X = \left\{ x \in M \mid \phi[x, b_1, \dots, b_n] \right\} \]
for some (fixed) choice of parameters $b_1,\dots,b_n \in M$.
We allow $n=0$, in which case we say $X$ is \vocab{$M$-definable without parameters}.
Note that $X$ need not itself be in $M$!
As a trivial example, $X = M$ is $M$-definable without parameters
(just take $\phi[x]$ to always be true), and certainly we do not have $X \in M$.
\begin{exercise}
	Verify (i)-(iv) above.
	(The others are not much harder, but more annoying.)
\end{exercise}
\begin{remark}
	Converses to the statements of \Cref{lem:transitive_ZFC}
	are true for all claims other than (vii).
\end{remark}

\section{Mostowski Collapse}
Up until now I have been only talking about transitive models,
because they were easier to think about.
Here's a second, better reason we might only care about transitive models.

\begin{lemma}
	[Mostowski Collapse]
	Let $\mathscr X = (X,E)$ be a model
	such that $\mathscr X \vDash \Extensionality + \Foundation$.
	Then there exists an isomorphism $\pi : \mathscr X \to M$ for
	a transitive model $M = (M,\in)$.
\end{lemma}

This is also called the \emph{transitive collapse}.
In fact, both $\pi$ and $M$ are unique.

\begin{proof}
	The idea behind the proof is very simple.
	Since $E$ is well-founded and extensional, we can look at the
	$E$-minimal element $x_\varnothing$ of $X$ with respect to $E$.
	Clearly, we want to send that to $0 = \varnothing$.

	Then we take the next-smallest set under $E$, and send it to $1 = \{\varnothing\}$.
	We ``keep doing this''; it's not hard to see this does exactly what we want.

	To formalize, define $\pi$ by transfinite recursion:
	\[ \pi(x) \defeq \left\{ \pi(y) \mid y \; E \; x \right\}. \]
	This $\pi$, by construction, does the trick.
\end{proof}

The picture of this is quite ``collapsing'' the elements of $M$ down
to the bottom of $V$, hence the name.
\missingfigure{Picture of Mostowski collapse}


\section{Adding an Inaccessible, Skolem Hulls, and Going Insane}
\prototype{$V_\kappa$}
At this point you might be asking, well, where's my model of $\ZFC$?

I unfortunately have to admit now: $\ZFC$ can never prove that there is a model of $\ZFC$
(unless $\ZFC$ is inconsistent, but that would be even worse).
This is a result called G\"odel's Incompleteness Theorem.

Nonetheless, with some very modest assumptions added, we can actually show that a model \emph{does} exist:
for example, assuming that there exists a strongly inaccessible cardinal $\kappa$ would do the trick,
$V_\kappa$ will be such a model (\Cref{prob:inaccessible_model}).
Intuitively you can see why: $\kappa$ is so big that any set of rank lower than it can't escape it
even if we take their power sets, or any other method that $\ZFC$ lets us do.
I encourage you to try \Cref{prob:inaccessible_model} now.

More pessimistically, this shows that it's impossible to prove in $\ZFC$ that such a $\kappa$ exists.
Nonetheless, we now proceed under $\ZFC^+$ for convenience, which adds the existence of such a $\kappa$
as a final axiom.
So we now have a model $V_\kappa$ to play with. Joy!

Great. Now we do something \emph{really} crazy.
\begin{theorem}[Countable Transitive Model]
	Assume $\ZFC^+$. Then there exists a transitive model $M$ of $\ZFC$
	such that $M$ is a \emph{countable} set.
\end{theorem}
\begin{proof}
	I'm telling you, this is insane. Fasten your seat belts.

	Start with the set $X_0 = \varnothing$.
	Then for every integer $n$, we do the following to get $X_{n+1}$.
	\begin{itemize}
		\ii Start with $X_{n+1}$ containing very element of $X_n$.
		\ii Consider a formula $\phi(x, x_1, \dots, x_n)$
		and $b_1, \dots, b_n$ in $X_n$.
		Suppose that $M$ thinks there is an $b \in M$ for which
		\[ M \vDash \phi[b, b_1, \dots, b_n]. \]
		We then add in the element $b$ to $X_{n+1}$.
		\ii We do this for \emph{EVERY POSSIBLE FORMULA IN THE LANGUAGE OF SET THEORY}.
		We also have to put in \emph{EVERY POSSIBLE SET OF PARAMETERS} from the previous set $X_n$.
	\end{itemize}
	At every step $X_n$ is countable.
	Reason: there are countably many possible finite sets of parameters in $X_n$,
	and countably many possible formulas, so in total we only ever add in countably many things
	at each step.
	This exhibits an infinite nested sequence of countable sets
	\[ X_0 \subseteq X_1 \subseteq X_2 \subseteq \dots \]
	None of these is a substructure of $M$, because each $X_n$ by relies on witnesses in $X_{n+1}$.
	So we instead \emph{TAKE THE UNION}:
	\[ X = \bigcup_n X_n. \]
	This satisfies the Tarski-Vaught test, and is countable.

	There is one minor caveat: $X$ might not be transitive.
	We don't care, because we just take its Mostowski collapse.
\end{proof}

Please take a moment to admire how insane this is.
It hinges irrevocably on the fact that there are countably many sentences we can write down.

\begin{remark}
	This proof relies heavily on the Axiom of Choice
	when we add in the element $b$ to $X_{n+1}$.
	Without Choice, there is no way of making these decisions all at once.

	Usually, the right way to formalize the Axiom of Choice usage is,
	for every formula $\phi(x, x_1, \dots, x_n)$, to pre-commit (at the very beginning)
	to a function $f_\phi(x_1, \dots, x_n)$, such that given any $b_1, \dots, b_n$
	$f_\phi(b_1, \dots, b_n)$ will spit out the suitable value of $b$ (if one exists).
	Personally, I think this is hiding the spirit of the proof, but it does
	make it clear how exactly Choice is being used.

	These $f_\phi$'s have a name: \vocab{Skolem functions}.
\end{remark}

The trick we used in the proof works in more general settings:
\begin{theorem}
	[Downward L\"owenheim-Skolem Theorem]
	Let $\MM = (M,E)$ be a model, and $A \subseteq M$.
	Then there exists a set $B$ (called the \vocab{Skolem hull} of $A$)
	with $A \subseteq B \subseteq M$,
	such that $(B,E) \prec \MM$, and
	\[ \left\lvert B \right\rvert < \max \left\{ \omega, \left\lvert A \right\rvert \right\}. \]
\end{theorem}
In our case, what we did was simply take $A$ to be the empty set.
\begin{ques}
	Prove this. (Exactly the same proof as before.)
\end{ques}


\section{FAQ's on Countable Models}
The most common one is ``how is this possible?'',
with runner-up ``what just happened''.

Let me do my best to answer the first question.
It seems like there are two things running up against each other:
\begin{enumerate}[(1)]
	\ii $M$ is a transitive model of $\ZFC$, but its universe is uncountable.
	\ii $\ZFC$ tells us there are uncountable sets!
\end{enumerate}
(This has confused so many people it has a name, Skolem's paradox.)

The reason this works I actually pointed out earlier:
\emph{countability is not absolute, it is $\Sigma_1$}.

Recall that a set $x$ is countable if
\emph{there exists} an injective map $x \injto \omega$.
The first statement just says that \emph{in the universe $V$},
there is a injective map $F: M \injto \omega$.
In particular, for any $x \in M$ (hence $x \subseteq M$, since $M$ is transitive),
$x$ is countable \emph{in $V$}.
This is the content of the first statement.

But for $M$ to be a model of $\ZFC$, $M$ only has to think statements in $\ZFC$ are true.
More to the point, the fact that $\ZFC$ tells us there are uncountable sets means
\[ M \vDash \text{$\exists x$ uncountable}. \]
In other words,
\[ M \vDash \exists x \forall f
	\text{ If $f : x \to \omega$ then $f$ isn't injective}. \]
The key point is the $\forall f$ searches only functions in our tiny model $M$.
It is true that in the ``real world'' $V$, there are injective functions $f : x \to \omega$.
But $M$ has no idea they exist!
It is a brain in a vat: $M$ is oblivious to any information outside it.

So in fact, every ordinal which appears in $M$ is countable in the real world.
It is just not countable in $M$.
Since $M \vDash \ZFC$, $M$ is going to think there is some smallest uncountable cardinal,
say $\aleph_1^M$.
It will be the smallest (infinite) ordinal in $M$
with the property that there is no bijection \emph{in the model $M$}
between $\aleph_1^M$ and $\omega$.
However, we necessarily know that such a bijection is going to exist in the real world $V$.

Put another way, cardinalities in $M$ can look vastly different from those in the real world,
because cardinality is measured by bijections, which I guess is inevitable, but leads to chaos.

\section{Picturing Inner Models}



\section\problemhead
\begin{sproblem}
	Show that for any transitive model $M$, the set of ordinals in $M$
	is itself some ordinal.
\end{sproblem}

\begin{dproblem}
	Assume $\MM_1 \subseteq \MM_2$. Show that
	\begin{enumerate}[(a)]
		\ii If $\phi$ is $\Delta_0$,
		then $\MM_1 \vDash \phi[b_1, \dots, b_n] \iff \MM_2 \vDash \phi[b_1, \dots, b_n]$.
		\ii If $\phi$ is $\Sigma_1$,
		then $\MM_1 \vDash \phi[b_1, \dots, b_n] \implies \MM_2 \vDash \phi[b_1, \dots, b_n]$.
		\ii If $\phi$ is $\Pi_1$,
		then $\MM_2 \vDash \phi[b_1, \dots, b_n] \implies \MM_1 \vDash \phi[b_1, \dots, b_n]$.
	\end{enumerate}
	(This should be easy if you've understood the chapter.)
\end{dproblem}

\begin{dproblem}
	[Tarski-Vaught Test]
	<++>
\end{dproblem}

\begin{sproblem}
	[Inaccessible Cardinal]
	\label{prob:inaccessible_model}
\end{sproblem}<++>

Reflection


\chapter{Forcing}
We are now going to introduce Paul Cohen's technique of \vocab{forcing},
which we then use to break the Continuum Hypothesis.

Here is how it works.
Given a transitive model $M$ and a poset $\Po$ inside it,
we can consider a ``generic'' subset $G \subseteq \Po$, where $G$ is not in $M$.
Then, we are going to construct a bigger universe $M[G]$ which contains both $M$ and $G$.
(This notation is deliberately the same as $\ZZ[\sqrt2]$, for example -- in the algebra case,
we are taking $\ZZ$ and adding in a new element $\sqrt 2$, plus everything that can be generated from it.)
By choosing $\Po$ well, we can cause $M[G]$ to have desirable properties.

The models $M$ and $M[G]$ will share the same ordinals.
But one issue with this is that forcing may introduce some new bijections between cardinals of $M$
that were not there originally; this leads to a phenomenon called \emph{cardinal collapse}:
quite literally, cardinals in $M$ will no longer be cardinals in $M[G]$, and instead just an ordinal.
\missingfigure{universe showing cardinal collapse}

In the case of the Continuum Hypothesis, we'll introduce a $\Po$ such that
any generic subset $G$ will ``encode'' $\aleph_2^M$ real numbers.
We'll then show cardinal collapse does not occur, meaning $\aleph_2^{M[G]} = \aleph_2^M$.
Thus $M[G]$ will have $\aleph_2^{M[G]}$ real numbers, as desired.

\section{Setting Up Posets}
\prototype{Infinite Binary Tree}
Let $M$ be a transitive model of $\ZFC$.
Let $\Po = (\Po, \le) \in M$ be a poset with a maximal element $1_\Po$
which lives inside a model $M$.
The elements of $\Po$ are called \vocab{conditions};
because they will force things to be true in $M[G]$.

\begin{definition}
	A subset $D \subseteq \Po$ is \vocab{dense} if for all $p \in \Po$,
	there exists a $q  \in D$ such that $q \le p$.
\end{definition}
Examples of dense subsets include the entire $\Po$ as well
as any downwards ``slice''.

\begin{definition}
	For $p,q \in \Po$ we write $p \parallel q$,
	saying ``$p$ is \vocab{compatible} with $q$'',
	if there is exists $r \in \Po$ with $r \le p$ and $r \le q$.
	Otherwise, we say $p$ and $q$ are \vocab{incompatible}
	and write $p \perp q$.
\end{definition}
\begin{example}[Infinite Binary Tree]
	Let $\Po = 2^{<\omega}$ be the \vocab{infinite binary tree} shown below,
	extended to infinity in the obvious way:
	\begin{center}
		\begin{asy}
			size(8cm);
			pair P = Drawing("\varnothing", (0,4), dir(90));
			pair P0 = Drawing("0", (-5,2), 1.5*dir(90));
			pair P1 = Drawing("1", (5,2),  1.5*dir(90));
			pair P00 = Drawing("00", (-7,0), 1.4*dir(120));
			pair P01 = Drawing("01", (-3,0), 1.4*dir(60));
			pair P10 = Drawing("10", (3,0),  1.4*dir(120));
			pair P11 = Drawing("11", (7,0),  1.4*dir(60));

			pair P000 = Drawing("000", (-8,-3));
			pair P001 = Drawing("001", (-6,-3));
			pair P010 = Drawing("010", (-4,-3));
			pair P011 = Drawing("011", (-2,-3));

			pair P100 = Drawing("100", (2,-3));
			pair P101 = Drawing("101", (4,-3));
			pair P110 = Drawing("110", (6,-3));
			pair P111 = Drawing("111", (8,-3));

			label("$\vdots$", (-7,-3), dir(-90));
			label("$\vdots$", (-3,-3), dir(-90));
			label("$\vdots$", (3,-3), dir(-90));
			label("$\vdots$", (7,-3), dir(-90));

			draw(P01--P0--P00);
			draw(P11--P1--P10);
			draw(P0--P--P1);
			draw(P000--P00--P001);
			draw(P100--P10--P101);
			draw(P010--P01--P011);
			draw(P110--P11--P111);
		\end{asy}
	\end{center}

	\begin{enumerate}[(a)]
		\ii The maximal element $1_\Po$ is the empty string $\varnothing$.
		\ii $D = \{\text{all strings ending in $001$}\}$ is an example of a dense set.
		\ii No two elements of $\Po$ are compatible unless they are comparable.
	\end{enumerate}
\end{example}


Now, I can specify what it means to be ``generic''.
\begin{definition}
	A nonempty set $G \subseteq \Po$ is a \vocab{filter} if
	\begin{enumerate}[(a)]
		\ii The set $G$ is upwards-closed:
		$\forall p \in G (\forall q \ge p) (q \in G)$.
		\ii Any pair of elements in $G$ is compatible.
	\end{enumerate}
	We say $G$ is \vocab{$M$-generic} if for all $D$ which are \emph{in the model $M$},
	if $D$ is dense then $G \cap D \neq \varnothing$.
\end{definition}
\begin{ques}
	Show that if $G$ is a filter then $1_\Po \in G$.
\end{ques}
\begin{example}[Generic Filters on the Infinite Binary Tree]
	Let $\Po = 2^{<\omega}$.
	The generic filters on $\Po$ are sets of the following form:
	\[ \left\{ 0,\; b_1,\; b_1b_2,\; b_1b_2b_3,\; \dots \right\}. \]
	So every generic filter on $\Po$ correspond to a binary number $b = 0.b_1b_2b_3\dots$.

	It is harder to describe which reals correspond to generic filters,
	but they should really ``look random''.
	For example, the set of strings ending in $011$ is dense,
	so one should expect ``$011$'' to appear inside $b$,
	and more generally that $b$ should contain every binary string.
	So one would expect the binary expansion of $\pi-3$ might correspond to a generic,
	but not something like $0.010101\dots$.
	That's why we call them ``generic''.
\end{example}

\begin{center}
	\begin{asy}
		size(8cm);
		pair P = Drawing("\varnothing", red, (0,4), red, dir(90));
		pair P0 = Drawing("0", red, (-5,2), red, 1.5*dir(90));
		pair P1 = Drawing("1", (5,2),  1.5*dir(90));
		pair P00 = Drawing("00", (-7,0), 1.4*dir(120));
		pair P01 = Drawing("01", red, (-3,0), red, 1.4*dir(60));
		pair P10 = Drawing("10", (3,0),  1.4*dir(120));
		pair P11 = Drawing("11", (7,0),  1.4*dir(60));

		pair P000 = Drawing("000", (-8,-3));
		pair P001 = Drawing("001", (-6,-3));
		pair P010 = Drawing("010", red, (-4,-3), red);
		pair P011 = Drawing("011", (-2,-3));

		pair P100 = Drawing("100", (2,-3));
		pair P101 = Drawing("101", (4,-3));
		pair P110 = Drawing("110", (6,-3));
		pair P111 = Drawing("111", (8,-3));

		draw(P01--P0--P00);
		draw(P11--P1--P10);
		draw(P0--P--P1);
		draw(P000--P00--P001);
		draw(P100--P10--P101);
		draw(P010--P01--P011);
		draw(P110--P11--P111);

		draw(P--P0--P01--P010--(P010+2*dir(-90)), red+1.4);
		MP("G", P010+2*dir(-90), dir(-90), red);
	\end{asy}
\end{center}

\begin{exercise}
	Verify that these are every generic filter $2^{<\omega}$ has the form above.
	Show that conversely, a binary number gives a filter, but it need not be generic.
\end{exercise}

Notice that if $p \ge q$, then the sentence $q \in G$ tells us more information than the sentence $p \in G$.
In that sense $q$ is a \emph{stronger} condition.
In another sense $1_\Po$ is the weakest possible condition,
because it tells us nothing about $G$; we always have $1_\Po \in G$
since $G$ is upwards closed.

\section{More Properties of Posets}
We had better make sure that generic filters exist.
In fact this is kind of tricky, but for countable models it works:
\begin{lemma}[Rasiowa-Sikorski Lemma]
	Suppose $M$ is a \emph{countable} transitive model of $\ZFC$
	and $\Po$ is a partial order.
	Then there exists an $M$-generic filter $G$.
\end{lemma}
\begin{proof}
	Since $M$ is countable, there are only countably many dense sets (they live in $M$!),
	say \[ D_1, D_2, \ldots, D_n, \ldots \in M. \]
	Using Choice,
	let $p_1 \in D_1$, and then let $p_2 \le p_1$ such that $p_2 \in D_2$
	(this is possible since $D_2$ is dense), and so on.
	In this way we can inductively exhibit a chain
	\[ p_1 \ge p_2 \ge p_3 \ge \dots \]
	with $p_i \in D_i$ for every $i$.

	Hence, we want to generate a filter from the $\{p_i\}$.
	Just take the upwards closure -- let $G$ be the set of $q \in \Po$ such that $q \ge p_n$ for some $n$.
	By construction, $G$ is a filter (this is actually trivial).
	Moreover, $G$ intersects all the dense sets by construction.
	\todo{problem?}
\end{proof}
Fortunately, for breaking $\CH$ we would want $M$ to be countable anyways.
% This is really just the proof of the Baire category theorem.

The other thing we want to do to make sure we're on the right track is guarantee
that a generic set $G$ is not actually in $M$.
(Analogy: $\ZZ[3]$ is a really stupid extension.)
The condition that guarantees this is:

\begin{definition}
	A partial order $\Po$ is \vocab{splitting} if
	for all $p \in \Po$, there exists $q,r \le p$
	such that $q \perp r$.
\end{definition}
\begin{example}[Infinite Binary Tree is (Very) Splitting]
	The infinite binary tree is about as splitting as you can get.
	Given $p \in 2^{<\omega}$, just consider the two elements right under it.
\end{example}

\begin{lemma}[Splitting Posets Omit Generic Sets]
	Suppose $\Po$ is splitting.  Then if $F \subseteq \Po$ is a filter
	such that $F \in M$, then $\Po \setminus F$ is dense.
	In particular, if $G \subseteq \Po$ is generic, then $G \notin M$.
\end{lemma}
\begin{proof}
	Consider $p \notin \Po \setminus F \iff p \in F$.
	Then there exists $q, r \le p$ which are not compatible.
	Since $F$ is a filter it cannot contain both;
	we must have one of them outside $F$, say $q$.
	Hence every element of $p \in \Po \setminus (\Po \setminus F)$
	has an element $q \le p$ in $\Po \setminus F$.
	That's enough to prove $\Po \setminus F$ is dense.
	\begin{ques}
		Deduce the last assertion of the lemma about generic $G$. \qedhere
	\end{ques}
\end{proof}

\section{Names, and the Generic Extension}
We now define the \emph{names} associated to a poset $\Po$.

\begin{definition}
	Suppose $M$ is a transitive model of $\ZFC$, $\Po = (\Po, \le) \in M$ is a partial order.
	We define the hierarchy of \vocab{$\Po$-names} recursively by
	\begin{align*}
		\Name_0 &= \varnothing \\
		\Name_{\alpha+1} &= \PP(\Name_\alpha \times \Po) \\
		\Name_{\lambda} &= \bigcup_{\alpha < \lambda} \Name_\alpha.
	\end{align*}
	Finally, $\Name = \bigcup_\alpha \Name_\alpha$ denote the class of all $\Po$-names.
	% For $\tau \in \Name$, let $\nrank(\tau)$ be the least $\alpha$ such that $\tau \in \Name_\alpha$.
\end{definition}
(These $\Name_\alpha$'s are the analog of the $V_\alpha$'s:
each $\Name_\alpha$ is just the set of all names with rank $\le \alpha$.)

\begin{definition}
	For a filter $G$, we define the \vocab{interpretation} of $\tau$ by $G$,
	denote $\tau^G$, using the transfinite recursion
	\[ \tau^G
		= \left\{ \sigma^G
		\mid \left<\sigma, p\right> \in \tau
		\text{ and } p \in G\right\}. \]
	We then define the model
	\[ M[G] = \left\{ \tau^G \mid \tau \in \Name^M \right\}. \]
	In words, $M[G]$ is the interpretation of all the possible $\Po$-names
	(as computed by $M$).
\end{definition}

\textbf{You should think of a $\Po$-name as a ``fuzzy set''.}
Here's the idea.
Ordinary sets are collections of ordinary sets,
so fuzzy sets should be collections of fuzzy sets.
These fuzzy sets can be thought of like the Ghosts of Christmases yet to come:
they represent things that might be, rather than things that are certain.
In other words, they represent the possible futures of $M[G]$ for various choices of $G$.

Every fuzzy set has an element $p \in \Po$ pinned to it.
When it comes time to pass judgment,
we pick a generic $G$ and filter through the universe of $\Po$-names.
The fuzzy sets with an element of $G$ attached to it materialize into the real world,
while the fuzzy sets with elements outside of $G$ fade from existence.
The result is $M[G]$.

\begin{example}[First Few Levels of the Name Hierarchy]
	Let us compute
	\begin{align*}
		\Name_0 &= \varnothing \\
		\Name_1 &= \PP(\varnothing \times \Po) \\
		&= \{\varnothing\} \\
		\Name_2 &= \PP(\{\varnothing\} \times \Po) \\
		&= \PP\left( \left\{ 
			\left<\varnothing, p\right>
			\mid p \in \Po
		\right\} \right).
	\end{align*}
\end{example}
Compare the corresponding von Neuman universe.
\[ V_0 = \varnothing, \; V_1 = \{\varnothing\}, \;
V_2 = \left\{ \varnothing, \left\{ \varnothing \right\} \right\}. \]

\begin{example}[Example of an Interpretation]
	As we said earlier, $\Name_1 = \{\varnothing\}$.
	Now suppose
	\[ \tau =
		\left\{
			\left<\varnothing, p_1\right>,
			\left<\varnothing, p_2\right>,
			\dots
			\left<\varnothing, p_n\right>
		\right\} 
		\in \Name_2. \]
	Then 
	\[
		\tau^G
		= \left\{ \varnothing \mid
		\left<\varnothing, p\right> \in \tau \text{ and } p \in G\right\}
		=
		\begin{cases}
			\{\varnothing\} & \text{if some } p_i \in G \\
			\varnothing & \text{otherwise}.
		\end{cases}
	\]
	In particular, remembering that $G$ is nonempty we see that
	\[ \left\{ \tau^G \mid \tau \in \Name_2 \right\} = V_2^M. \]
	In fact, this holds for any natural number $n$, not just $2$.
\end{example}
So, $M[G]$ and $M$ agree on finite sets.

Now, we want to make sure $M[G]$ contains the elements of $M$.
To do this, we take advantage of the fact that $1_\Po$ must be in $G$, and define
for every $x \in M$ the set
\[ \check x = \left\{ \left<\check y, 1_\Po\right> \mid y \in x \right\} \]
by transfinite recursion.
Basically, $\check x$ is just a copy of $x$ where we add check marks and tag every element with $1_\Po$.

\begin{example}
	Compute $\check 0 = 0$ and $\check 1 = \left\{ \left<\check 0, 1_\Po\right> \right\}$.
	Thus \[ (\check 0)^G = 0 \quad\text{and}\quad (\check 1)^G = 1. \]
\end{example}
\begin{ques}
	Show that in general, $(\check x)^G = x$.
	(Rank induction.)
\end{ques}

However, we'd also like to cause $G$ to be in $M[G]$.
In fact, we can write down the name exactly:
\[ \dot G = \left\{ \left<\check p, p\right> \mid p \in \Po \right\}. \]
\begin{ques}
	Show that $(\dot G)^G = G$.
\end{ques}
\begin{ques}
	Verify that $M[G]$ is transitive:
	that is, if $\sigma^G \in \tau^G \in M[G]$, show that $\sigma^G \in M[G]$.
	(This is offensively easy.)
\end{ques}

In summary,
\begin{moral}
	$M[G]$ is a transitive model extending $M$ (it contains $G$).
\end{moral}

Moreover, it is reasonably well-behaved even if $G$ is just a filter.
Let's see what we can get off the bat.
\begin{lemma}[Properties Obtained from Filters]
	Let $M$ be a transitive model of $\ZFC$.
	If $G$ is a filter, then $M[G]$ is transitive
	and satisfies Extensionality, Foundation, EmptySet, Infinity, Pairing, and Union.
\end{lemma}
\todo{mathtt these}
This leaves PowerSet, Comprehension, Replacement, and Choice.
\begin{proof}
	Hence, we get Extensionality and Foundation for free.
	Infinity and EmptySet follows from $M \subseteq M[G]$.

	For Pairing, suppose $\sigma_1^G, \sigma_2^G \in M[G]$.
	Then
	\[ \sigma = 
		\left\{ \left<\sigma_1, 1_\Po\right>, \left<\sigma_2, 1_\Po\right> \right\}
	\]
	works satisfies $\sigma^G = \{\sigma_1, \sigma_2\}$.
	(Note that we used Pairing in $M$.)
	Union is left as a problem, which you are encouraged to try now.
\end{proof}
Up to here, we don't need to know anything about when a sentence is true in $M[G]$;
all we had to do was contrive some names like $\check x$ or
$\left\{ \left<\sigma_1, 1_\Po\right>, \left<\sigma_2, 1_\Po\right> \right\}$
to get the facts we wanted.
But for the remaining axioms, we \emph{are} going to need this extra power
are true in $M[G]$.
For this, we have to introduce the Fundamental Theorem of Forcing.

\section{Fundamental Theorem of Forcing}
The model $M$ unfortunately has no idea what $G$ might be,
only that it is some generic filter.\footnote{%
	You might say this is a good thing, here's why.
	We're trying to show that $\neg \CH$ is consistent with $\ZFC$,
	and we've started with a model $M$ of the real universe $V$.
	But for all we know $\CH$ might be true in $V$ (what if $V=L$?),
	in which case it would also be true of $M$.

	Nonetheless, we boldly construct $M[G]$ an extension of the model $M$.
	In order for it to behave differently from $M$, it has to be out of reach of $M$.
	Conversely, if $M$ could compute everything about $M[G]$,
	then $M[G]$ would have to conform to $M$'s beliefs.

	That's why we worked so hard to make sure $G \in M[G]$ but $G \notin M$.
}
Nonetheless, we are going to define a relation $\Vdash$, called the \emph{forcing} relation.
Roughly, we are going to write
\[ p \Vdash \varphi(\sigma_1, \dots, \sigma_n) \]
where $p \in \Po$, $\sigma_1, \dots, \sigma_n \in M[G]$, if and only if the following holds:
\begin{quote}
	For \emph{any} generic $G$,
	if $p \in G$,
	then $M[G] \vDash \varphi(\sigma_1^G, \dots, \sigma_n^G)$.
\end{quote}
Note that $\Vdash$ is defined without reference to $G$:
it is something that $M$ can see.
We say $p$ \vocab{forces} the sentences $\varphi(\sigma_1, \dots, \sigma_n)$.
And miraculously, we can define this relation in such a way that the converse is true:
\emph{a sentence holds if and only if some $p$ forces it}.


\begin{theorem}
	[Fundamental Theorem of Forcing]
	Suppose $M$ is a transitive model of ZF.
	Let $\Po \in M$ be a poset, and $G \subseteq \PP$ is an $M$-generic filter.
	Then,
	\begin{enumerate}[(1)]
		\ii Consider $\sigma_1, \dots, \sigma_n \in \Name^M$,
		Then
		\[ M[G] \vDash \varphi[\sigma_1^G, \dots, \sigma_n^G] \]
		if and only if there exists a condition $p \in G$
		such that $p$ \emph{forces} the sentence $\varphi(\sigma_1, \dots, \sigma_n)$.
		We denote this by $p \Vdash \varphi(\sigma_1, \dots, \sigma_n)$.
		\ii This forcing relation is (uniformly) definable in $M$.
	\end{enumerate}
\end{theorem}

I'll tell you how the definition works in the following.

\section{(Optional) Defining the Relation}
Here's how we're going to go.
We'll define the most generous condition possible such that
the forcing works in one direction ($p \vDash \varphi(\sigma_1, \dots, \sigma_n)$ means
$M[G] \vDash \varphi[\sigma_1^G, \dots, \sigma_n^G]$).
We will then cross our fingers that the converse also works.

We proceed by induction on the formula complexity.
It turns out in this case that the atomic formula (base cases)
are hardest and themselves require induction on ranks.

For some motivation, let's consider how we should define $p \Vdash \tau_1 \in \tau_2$ given that we've already defined $p \vDash \tau_1 = \tau_2$.
We need to ensure this holds iff
\[ \forall \text{$M$-generic $G$ with $p \in G$}:
	\ M[G] \vDash \tau_1^G \in \tau_2^G. \]
So it suffices to ensure that any generic $G \ni p$ hits a condition $q$ which forces $\tau_1^G$ to \emph{equal} a member $\tau^G$ of $\tau_2^G$.
In other words, we want to choose the definition of $p \Vdash \tau_1 \in \tau_2$ to hold if and only if
\[
	\left\{ q \in \Po
	\mid \exists \left<\tau, r\right> \in \tau_2 
	\left( q \le r \land q \Vdash(\tau=\tau_1) \right)
	\right\}
\]
is dense below in $p$.
In other words, if the set is dense, then the generic must hit $q$, so it must hit $r$, meaning that $\left<\tau_r\right> \in \tau_2$ will get interpreted such that $\tau^G \in \tau_2^G$, and moreover the $q \in G$ will force $\tau_1 = \tau$.

Now let's write down the definition\dots
In what follows, the $\Vdash$ omits the $M$ and $\Po$.
\begin{definition}
	Let $M$ be a countable transitive model of ZFC.
	Let $\Po \in M$ be a partial order.
	For $p \in \Po$ and $\varphi(\sigma_1, \dots, \sigma_n)$ a formula in LST, we write $\tau \Vdash \varphi(\sigma_1, \dots, \sigma_n)$ to mean the following, defined by induction on formula complexity plus rank.
	\begin{enumerate}[(1)]
		\ii $p \Vdash \tau_1 = \tau_2$ means
		\begin{enumerate}[(i)]
			\ii For all $\left<\sigma_1, q_1\right> \in \tau_1$ the set
			\[ D_{\sigma_1, q_1}
				\defeq
				\left\{ r \mid
				r \le q_1 \lthen \exists \left<\sigma_2, q_2\right> \in \tau_2 \left( r \le q_2 \land r \Vdash (\sigma_1 = \sigma_2) \right)\right\}.
			\]
			is dense in $p$.
			(This encodes ``$\tau_1 \subseteq \tau_2$''.)
			\ii For all $\left<\sigma_2, q_2\right> \in \tau_2$,
			the set $D_{\sigma_2, q_2}$ defined similarly is dense below $p$.
		\end{enumerate}
		\ii $p \Vdash \tau_1 \in \tau_2$ means
		\[
		\left\{ q \in \Po
		\mid \exists \left<\tau, r\right> \in \tau_2 
		\left( q \le r \land q \Vdash(\tau=\tau_1) \right)
		\right\} \]
		is dense below $p$.
		\ii $p \Vdash \varphi \land \psi$ means $p \Vdash \varphi$ and $p \Vdash \psi$.
		\ii $p \Vdash \neg \varphi$ means $\forall q \le p$, $q \not\Vdash \varphi$.
		\ii $p \Vdash \exists x \varphi(x, \sigma_1, \dots, \sigma_n)$ means that the set
		\[
			\left\{ q \mid \exists \tau
				\left( q \Vdash \varphi(\tau, \sigma_1, \dots, \sigma_n \right)
			\right\}
		\]
		is dense below $p$.
	\end{enumerate}
\end{definition}
This is definable in $M$!
All we've referred to is $\Po$ and names, which are in $M$.
(Note that being dense is definable.)
Actually, in parts (3) through (5) of the definition above,
we use induction on formula complexity.
But in the atomic cases (1) and (2) we are doing induction on the ranks of the names.

So, the construction above gives us one direction (I've omitted tons of details, but\dots).

Now, how do we get the converse: that a sentence is true if and only if something forces it?
Well, by induction, we can actually show the following result:
\begin{lemma}[Consistency and Persistence]
	We have
	\begin{enumerate}[(1)]
		\ii (Consistency) If $p \Vdash \varphi$ and $q \le p$ then $q \Vdash \varphi$.
		\ii (Persistence) If $\left\{ q \mid q \Vdash \varphi \right\}$
		is dense below $p$ then $p \Vdash \varphi$.
	\end{enumerate}
\end{lemma}
These are just inductions on the five parts of the definition.
From this it also follows that
\begin{corollary}[Completeness]
	The set $\left\{ p \mid p \Vdash \varphi \text{ or } p \Vdash \neg\varphi \right\}$
	is dense.
\end{corollary}
\begin{proof}
	We claim that whenever $p \not\Vdash \varphi$ then
	for some $\ol p \le p$ we have $\ol p \Vdash \neg\varphi$;
	this will establish the corollary.

	By the contrapositive of the previous lemma,
	$\{q \mid q \Vdash \varphi\}$ is not dense below $p$,
	meaning for some $\ol p \le p$, every $q \le \ol p$ gives $q \not\Vdash \varphi$.
	By the definition of $p \vDash \neg\varphi$,
	we have $\ol p \vDash \neg\varphi$.
\end{proof}
And this gives the converse: the $M$-generic $G$ has to hit some condition
that passes judgment, one way or the other.
This completes the proof of the Fundamental Theorem.

\section{The Remaining Axioms}
\begin{theorem}[The Generic Extension Satisfies $\ZFC$]
	Suppose $M$ is a transitive model of $\ZFC$.
	Let $\Po \in M$ be a poset, and $G \subseteq \PP$ is an $M$-generic filter.
	Then \[ M[G] \vDash \ZFC. \]
\end{theorem}
\begin{proof}
	We'll just do Comprehension, as the other remaining axioms are similar.
	
	Suppose $\sigma^G, \sigma_1^G, \dots, \sigma_n^G \in M[G]$
	are a set and parameters, and
	$\varphi(x,x_1, \dots, x_n)$ is an LST formula.
	We want to show that the set
	\[ A = \left\{ 
		x \in \sigma^G \mid M[G] \vDash \varphi[x, \sigma_1^G, \dots, \sigma_n^G]
	\right\} \]
	is in $M[G]$; i.e.\ it is the interpretation of some name.

	Note that every element of $\sigma^G$ is of the form $\rho^G$
	for some $\rho \in \dom(\sigma)$ (a bit of abuse here,
	$\sigma$ is a bunch of pairs of names and $p$'s,
	and the domain is just the set of names).
	So by the Fundamental Theorem of Forcing, we may write
	\[ A = 
		\left\{ \rho^G \mid \rho \in \dom(\sigma)
			\text{ and }
			\exists p \in G
			\left( p \Vdash \rho \in \sigma
			\land \varphi(\rho, \sigma_1, \dots, \sigma_n)
			\right)
		\right\}.
	\]
	To show $A \in M[G]$ we have to write down a $\tau$
	such that the name $\tau^G$ coincides with $A$.
	We claim that
	\[
		\tau
		=
		\left\{ \left<\rho, p\right>
			\in \dom(\sigma) \times \Po \mid
			p \Vdash 
			\land \varphi(\rho, \sigma_1, \dots, \sigma_n)
		\right\}
	\]
	is the correct choice.
	It's actually clear that $\tau^G = A$ by construction;
	the ``content'' is showing that $\tau$ is in actually a name of $M$,
	which follows from Comprehension in $M$.

	So really, the point of the Fundamental Theorem of Forcing
	is just to let us write down this $\tau$;
	it lets us show that $\tau$ is in $\Name^M$
	without actually referencing $G$.
\end{proof}


\section\problemhead
RS Lemma

Union



%\begin{exercise}
%	Show that $\rank \sigma^G \le \nrank(\sigma)$ for any $\sigma \in \Name^M$.
%\end{exercise}

%\begin{exercise}
%	Check that
%	\begin{enumerate}[(1)]
%		\ii $(\check x)^G = x$.
%		\ii $(\dot G)^G = G$.
%	\end{enumerate}
%\end{exercise}



\chapter{The Failure of the Continuum Hypothesis}
We now use the technique of forcing to break the Contiuum Hypothesis by choosing a good poset $\Po$.

\section{Forcing $V \neq L$ is really easy}
As a small aside, to check we're on the right track we show the following result.

\begin{theorem}[$V \ne L$]
	Let $M$ be a countable transitive model of $\ZFC$.
	Let $\Po \in M$ be \emph{any} splitting poset,
	and let $G \subseteq \Po$ be $M$-generic.
	Then $M[G] \vDash (V \neq L)$.
\end{theorem}
\begin{proof}
	Since $L$ has a $\Sigma_1$ definition,
	we have \[ L^{M[G]} = L^M \subseteq M \subsetneq M[G] \]
	where the last part follows from $G \notin M[G]$.
\end{proof}

Thus $M[G] \vDash \ZFC + (V \ne L)$ for any splitting poset $\Po$,
and we are one step closer to breaking $\CH$.

\section{Adding in Reals}
Starting with a \emph{countable} transitive model $M$.

We want to choose $\Po \in M$ such that $(\aleph_2)^M$ many real numbers appear,
and then worry about cardinal collapse later.

Recall the earlier situation where we set $\Po$ to be the infinite complete binary tree; its nodes can be thought of as partial functions $n \to 2$ where $n < \omega$.
Then $G$ itself is a path down this tree; i.e.\ it can be encoded as a total function $G : \omega \to 2$,
and corresponds to a real number.

\begin{center}
	\begin{asy}
		size(8cm);
		pair P = Drawing("\varnothing", red, (0,4), red, dir(90));
		pair P0 = Drawing("0", red, (-5,2), red, 1.5*dir(90));
		pair P1 = Drawing("1", (5,2),  1.5*dir(90));
		pair P00 = Drawing("00", (-7,0), 1.4*dir(120));
		pair P01 = Drawing("01", red, (-3,0), red, 1.4*dir(60));
		pair P10 = Drawing("10", (3,0),  1.4*dir(120));
		pair P11 = Drawing("11", (7,0),  1.4*dir(60));

		pair P000 = Drawing("000", (-8,-3));
		pair P001 = Drawing("001", (-6,-3));
		pair P010 = Drawing("010", red, (-4,-3), red);
		pair P011 = Drawing("011", (-2,-3));

		pair P100 = Drawing("100", (2,-3));
		pair P101 = Drawing("101", (4,-3));
		pair P110 = Drawing("110", (6,-3));
		pair P111 = Drawing("111", (8,-3));

		draw(P01--P0--P00);
		draw(P11--P1--P10);
		draw(P0--P--P1);
		draw(P000--P00--P001);
		draw(P100--P10--P101);
		draw(P010--P01--P011);
		draw(P110--P11--P111);

		draw(P--P0--P01--P010--(P010+2*dir(-90)), red+1.4);
		MP("G", P010+2*dir(-90), dir(-90), red);
	\end{asy}
\end{center}

We want to do something similar, but with $\omega_2$ many real numbers instead of just one.
In light of this, consider in $M$ the following poset:
\[
	\Po = 
	\opname{Add} \left( \omega_2, \omega \right)
	\defeq
	\left( 
	\left\{ p : \omega_2 \times \omega \to 2,
		\dom(p) < \omega
	\right\},
	\supseteq
	\right).
\]
These elements (conditions) are ``partial functions'':
we take some finite subset of $\omega \times \omega_2$ and map it into $2=\{0,1\}$.
Moreover, we say $p \le q$ if $\dom(p) \supseteq \dom(q)$ and the two functions agree over $\dom(q)$.
\begin{ques}
	What is $1_\Po$ here?
\end{ques}

\begin{exercise}
	Show that a generic $G$ can be encoded as a function $\omega_2 \times \omega \to 2$.
\end{exercise}

%Let $G \subseteq \opname{Add}(\omega_2, \omega)$ be an $M$-generic.
%We claim that, like in the binary case, $G$ can be encoded as a function $\omega_2 \times \omega \to 2$.
%To see this, consider $\alpha \in \omega_2$ and $n \in \omega$; we have the dense set
%\[ D_{\alpha, n}
%	= \left\{ p \in \opname{Add}(\omega_2, \omega)
%	\mid (\alpha, n) \in \dom(p) \right\}
%\]
%(this is obviously dense, given any $p$ add in $(\alpha, n)$ if it's not in there already).
%So $G$ hits this dense set, meaning that for every $(\alpha, n)$ there's a function in $G$ which defines it.
%Using the fact that $G$ is upwards closed and a filter, we may as before we may interpret $G$ as a function $\omega_2 \times \omega \to 2$.

\begin{lemma}[$G$ encodes distinct real numbers]
	For $\alpha \in \omega_2$ define
	\[ G_\alpha = \left\{ n \mid G\left( \alpha,n \right) = 0 \right\} \in \PP(\NN). \]
	Then $G_\alpha \neq G_\beta$ for any $\alpha \neq \beta$.
\end{lemma}
\begin{proof}
	We claim that the set
	\[ D = \left\{ q \mid \exists n \in \omega :
		q\left( \alpha, n \right) \neq q\left( \beta, n \right)
		\text{ are both defined}
	\right\} \]
	is dense.
	\begin{ques}
		Check this.
		(Use the fact that the domains are all finite.)
	\end{ques}
%	This is pretty easy to see.
%	Consider $p \in \opname{Add}(\omega_2, \omega)$.
%	Then you can find an $n$ such that
%	neither $(\alpha, n)$ nor $(\beta, n)$ is defined,
%	just because $\dom(p)$ is finite.
%	Then you make $p'$ as $p$ plus $p'( (\alpha, n) ) = 1$
%	and $p'( (\beta, n) ) = 0$.
%	Hence the set is dense.

	Since $G$ is an $M$-generic it hits this dense set $D$.
	Hence $G_\alpha \neq G_\beta$.
\end{proof}

Since $G \in M[G]$ and $M[G] \vDash \ZFC$,
it follows that each $G_\alpha$ is in $M[G]$.
So there are at least $\aleph_2^M$ real numbers in $M$.
We are done once we can show there is no cardinal collapse.

\section{The Countable Chain Condition}
It remains to show that with $\Po = \opname{Add}(\omega, \omega_2)$, we have that
\[ \aleph_2^{M[G]} = \aleph_2^M. \]
In that case, since $M[G]$ will have $\aleph_2^M = \aleph_2^{M[G]}$ many reals, we will be done.

To do this, we'll rely on the following combinatorial property of $\Po$:

\begin{definition}
	We say that $A \subset \mathcal P$ is a \vocab{strong antichain}
	if for any distinct $p$ and $q$ in $A$, we have $p \perp q$.
\end{definition}
\begin{example}[Example of an Antichain]
	In the infinite binary tree, 
	the set $A = \{00, 01, 10, 11\}$ is a strong antichain
	(in fact maximal by inclusion).
\end{example}
This is stronger than the notion of ``antichain'' than you might be used to!\footnote{%
	In the context of forcing, some authors use ``antichain'' to refer to ``strong antichain''.
	I think this is lame.}
We don't merely require that every two elements are incomparable,
but that they are in fact \emph{incompatible}.
\begin{ques}
	Draw a finite poset and an antichain of it which is not strong.
\end{ques}

\begin{definition}
	A poset $\Po$ has the \vocab{$\kappa$-chain condition}
	(where $\kappa$ is a cardinal) if all strong antichains
	in $\Po$ have size less than $\kappa$.
	The special case $\kappa = \omega$ is called the \vocab{countable chain condition}.
\end{definition}

We are going to show that if the poset has the $\kappa$-chain condition
then it preserves all cardinals greater than $\kappa$.
% or was it > \kappa?
In particular, the countable chain condition will show that $\Po$ preserves all the cardinals.
Then, we'll show that $\opname{Add}(\omega, \omega_2)$ does indeed have this property.
This will complete the proof.

We isolate the following lemma.
\begin{lemma}[Possible Values Argument]
	Suppose $M$ is a transitive model of $\ZFC$ and $\Po$ is a partial order
	such that $\Po$ has the $\kappa$-chain condition in $M$.
	Let $X,Y \in M$ and let $f: X \to Y$
	be some function in $M[G]$, but $f \notin M$.

	Then there exists a function $F \in M$, with $F: X \to \PP(Y)$ and such that
	for any $x \in X$,
	\[ f(x) \in F(x) \quad\text{and}\quad \left\lvert F(x) \right\rvert^M < \kappa. \]
\end{lemma}
What this is saying is that if $f$ is some new function that's generated,
$M$ is still able to pin down the values of $f$ to at most $\kappa$ many values.

\begin{proof}
	The idea behind the proof is easy: any possible value of $f$ gives us some condition in
	the poset $\Po$ which forces it.
	Since distinct values must have incompatible conditions,
	the $\kappa$-chain condition guarantees
	there are at most $\kappa$ such values.

	Here are the details.
	Let $\dot f$, $\check X$, $\check Y$ be names for $f$, $X$, $Y$.
	Start with a condition $p$ such that $p$ forces the sentence
	\[ \text{``$\dot f$ is a function from $\check X$ to $\check Y$''}. \]
	We'll work just below here.

	For each $x \in X$, we can consider (using the Axiom of Choice) a maximal antichain $A(x)$
	of incompatible conditions $q \le p$ which forces $f(x)$ to equal some value $y \in Y$.
	Then, we let $F(x)$ collect all the resulting $y$-values.
	These are all possible values, and there are less than $\kappa$ of them.
\end{proof}

\section{Preserving Cardinals}
\begin{definition}
	For $M$ a transitive model of ZFC and $\Po \in M$ a poset,
	we say $\Po$ \vocab{preserves cardinals} if
	$\forall G \subseteq \Po$ an $M$-generic,
	the model $M$ and $M[G]$ agree on the sentence ``$\kappa$ is a cardinal'' for every $\kappa$.

	In the same way we will talk about $\Po$ preserving cofinalities, et cetera.
\end{definition}

\todo{write out this exercise}
\begin{exercise}
	Let $M$ be a transitive model of ZFC.
	Let $\Po \in M$ be a poset.
	Show that the following are equivalent for each $\lambda$:
	\begin{enumerate}[(1)]
		\ii $\Po$ preserves cofinalities less than or equal to $\lambda$.
		\ii $\Po$ preserves regular cardinals less than or equal to $\lambda$.
	\end{enumerate}
	Moreover the same holds if we replace ``less than or equal to''
	by ``greater than or equal to''.
\end{exercise}
Thus, to show that $\Po$ preserves cardinality and cofinalities it suffices to show that $\Po$ preserves regularity.

\begin{theorem}
	Suppose $M$ is a transitive model of ZFC, and $\Po \in M$ is a poset.
	Suppose $M$ satisfies the sentence ``$\Po$ has the $\kappa$ chain condition and $\kappa$ is regular''.
	Then $\Po$ preserves cardinals and cofinalities greater than or equal to $\kappa$.
\end{theorem}
\begin{proof}
	It suffices to show that $\Po$ preserves regularity greater than or equal to $\kappa$.
	Consider $\lambda > \kappa$ which is regular in $M$,
	and suppose for contradiction that $\lambda$ is not regular in $M[G]$.
	That's the same as saying that there is a function $f \in M[G]$,
	$f : \ol \lambda \to \lambda$ cofinal, with $\ol \lambda < \lambda$.
	Then by the Possible Values Argument,
	there exists a function $F \in M$ from $\ol \lambda \to \PP(\lambda)$
	such that $f(\alpha) \in F(\alpha)$ and $\left\lvert F(\alpha) \right\rvert^M < \kappa$
	for every $\alpha$.

	Now we work in $M$ again.
	Note for each $\alpha \in \ol\lambda$,
	$F(\alpha)$ is bounded in $\lambda$ since $\lambda$ is regular in $M$ and
	greater than $\left\lvert F(\alpha) \right\rvert$.
	Now look at the function $\ol \lambda \to \lambda$ in $M$ by just
	\[ \alpha \mapsto \cup F(\alpha) < \lambda. \]
	This is cofinal in $M$, contradiction.
\end{proof}

\subsection{Infinite Combinatorics}
In particular, if $\Po$ has the countable chain condition then $\Po$ preserves all the cardinals (and cofinalities).
Therefore, it remains to show that $\opname{Add}(\omega, \omega_2)$ satisfies the countable chain condition.
And this is going to be infinite combinatorics.

\begin{definition}
	Suppose $C$ is an uncountable collection of finite sets.
	$C$ is a \textbf{$\Delta$-system} if there exists a \textbf{root} $R$
	with the condition that for any distinct $X$ and $Y$
	in $C$, we have $X \cap Y = R$.
\end{definition}

\begin{lemma}
	[$\Delta$-System Lemma] Suppose $C$ is an uncountable collection of finite sets.
	Then $\exists \ol C \subseteq C$ such that
	\begin{enumerate}[(1)]
		\ii $\ol C$ is uncountable.
		\ii $\ol C$ is a $\Delta$-system.
	\end{enumerate}
\end{lemma}
\begin{proof}
	There exists an integer $n$ such that $C$ has uncountably many guys of length $n$.
	So we can throw away all the other sets, and just assume that all sets in $C$ have size $n$.

	We now proceed by induction on $n$.
	The base case $n=1$ is trivial, since we can just take $R = \varnothing$.
	For the inductive step we consider two cases.

	First, assume there exists an $a \in C$ contained in uncountably many $F \in C$.
	Throw away all the other guys.
	Then we can just delete $a$, and apply the inductive hypothesis.

	Now assume that for every $a$, only countably many members of $C$ have $a$ in them.
	We claim we can even get a $\ol C$ with $R = \varnothing$.
	First, pick $F_0 \in C$.
	It's straightforward to construct an $F_1$ such that $F_1 \cap F_0 = \varnothing$.
	And we can just construct $F_2, F_3, \dots$
\end{proof}

\begin{lemma}
	For all $\kappa$, $\opname{Add}(\omega, \kappa)$ satisfies the countable chain condition.
\end{lemma}
\begin{proof}
	Assume not. Let
	\[ \left\{ p_\alpha : \alpha < \omega_1 \right\} \]
	be an antichain.  Let
	\[ C = \left\{ \dom(p_\alpha) : \alpha < \omega_1 \right\}. \]
	Let $\ol C \subseteq C$ be such that $\ol C$ is uncountable, and $\ol C$ is a $\Delta$-system which root $R$.
	Then let
	\[ B = \left\{ p_\alpha : \dom(p_\alpha) \in R \right\}. \]
	Each $p_\alpha \in B$ is a function $p_\alpha : R \to \{0,1\}$,
	so there are two that are the same.
\end{proof}
\todo{blah blah blah}

\end{document}
